\section{Đánh giá kết quả đạt được}

\subsection{Kết quả đạt được so với kế hoạch}
Trong thời hạn 5 tuần, nhóm đã hoàn thành 14 hạng mục bàn giao đã liệt kê ở Chương~2, mục~\ref{sec:deliverables-list}:

\begin{table}[htbp]
    \centering
    \small
    \caption{Đối chiếu kế hoạch và kết quả thực hiện}
    \begin{tabularx}{\textwidth}{|Y|p{2.2cm}|Y|}
        \hline
        \textbf{Hạng mục} & \textbf{Trạng thái} & \textbf{Ghi chú} \\
        \hline
        SRS đầy đủ (FR, NFR, RTM) & Hoàn thành & \\
        \hline
        Architecture Design Document & Hoàn thành & \\
        \hline
                ERD & Hoàn thành & Đủ 20 thực thể, khớp schema.sql thật \\
        \hline
        UML Class Diagram & Đang thực hiện & Bản nháp đầu còn thiếu 8/20 bảng (Ledger, SalesChannel, Stocktake...), cần vẽ lại theo Entity Catalog (mục 4.2.1) \\
        \hline
        Use Case Diagram & Đang thực hiện & Cần đối chiếu lại phân quyền quản lý sản phẩm cho khớp Bảng RBAC (mục 3.3.3.3) \\
        \hline
        Sequence Diagram (6 luồng nghiệp vụ) & Đã có mã nguồn & Đủ mã PlantUML cho 6 luồng, đã qua nhiều vòng review logic, chờ xuất hình chính thức \\
        \hline
        Đặc tả API (Swagger/Postman) & Đang thực hiện & Đã có khung OpenAPI mẫu (6/\textasciitilde30 endpoint tiêu biểu), cần mở rộng đủ \\
        \hline
        Script DDL SQL Server & Hoàn thành & Đã chạy thật trên SQL Server 2022, tạo đủ 20 bảng \\
        \hline
        Trigger/rule append-only cho Ledger & Hoàn thành & Đã kiểm chứng PASS thật (TC-02, TC-03) \\
        \hline
        Seed data \& script kiểm tra toàn vẹn & Hoàn thành & Đã chạy thật, 13/13 test case (TC-01→TC-13) PASS \\
        \hline
        Backend API (.NET) & Chưa thực hiện & Ngoài phạm vi bàn giao học kỳ này, dời sang giai đoạn tiếp theo \\
        \hline
        Frontend (Admin Dashboard, POS PWA) & Chưa thực hiện & Ngoài phạm vi bàn giao học kỳ này \\
        \hline
    \end{tabularx}
\end{table}

\textit{Ghi chú: 8/13 hạng mục đã Hoàn thành với bằng chứng chạy thật (Chương 4, mục~\ref{sec:test-results}); 3 hạng mục đang hoàn thiện — chi tiết xem mục 5.3.1; 2 hạng mục ngoài phạm vi theo mục~\ref{sec:scope-limitations}.}
\subsection{Đánh giá theo yêu cầu phi chức năng}
Vì chưa có Backend thật, các chỉ tiêu phi chức năng về hiệu năng thời gian thực (API Response Time, khả năng chịu tải) chưa thể đo lường trực tiếp. Nhóm đã kiểm chứng được ở mức thiết kế/dữ liệu:
\begin{itemize}
    \item Ràng buộc append-only (NFR-SEC-03): đã kiểm chứng thành công qua trigger SQL Server.
    \item Cô lập dữ liệu đa tenant (NFR-TENANT-01): đã kiểm chứng qua truy vấn chéo TenantId trên seed data.
    \item Công thức giá vốn bình quân: đã kiểm chứng đúng với kịch bản nhập hàng nhiều đợt trong seed data.
    \item Concurrency (NFR-PERF-02): mới mô phỏng bằng nhiều phiên SQL song song, \textbf{chưa đo lường được} ở quy mô 50 request đồng thời qua API thật — cần Backend hoàn chỉnh để kiểm thử đầy đủ.
\end{itemize}

\subsection{Bài học kinh nghiệm}
\begin{itemize}
    \item Việc hoàn thành ERD và Data Dictionary trước khi viết DDL giúp nhóm tránh phải làm lại schema nhiều lần.
    \item Quyết định dùng TRIGGER thay vì RULE để chặn UPDATE/DELETE (Chương~4) giúp việc kiểm thử ràng buộc append-only rõ ràng và dễ chứng minh hơn với GVHD.
    \item Thu hẹp phạm vi bàn giao ngay từ đầu (chỉ Thiết kế + Database) giúp nhóm tập trung làm sâu và đúng phần lõi khó nhất của đề tài (cơ chế chống bán vượt tồn, sổ cái append-only), thay vì dàn trải sang Backend/Frontend mà không đủ thời gian hoàn thiện.
\end{itemize}
